%! Author = Admin
%! Date = 12.12.2025
%========================================================
% Памятки по загруженным статьям (обзор литературы)
%========================================================

\section*{1) \textit{Trajectory planning and tracking for UAVs with deep reinforcement learning and adaptive nonlinear MPC} \href{https://www.sciencedirect.com/science/article/pii/S0957417425042642}{ref}}
\begin{itemize}
  \item \textbf{Какая задача решалась.}
  Предложить метод \emph{планирования и слежения} за траекторией БПЛА в сложной динамической среде в реальном времени, улучшив устойчивость/точность по сравнению с фиксированно настроенным MPC.

  \item \textbf{Использовалось ли ИИ.}
  Да: используется \textbf{deep reinforcement learning} для \emph{динамической настройки} (адаптации) весов в целевой функции \textbf{adaptive nonlinear MPC (ANMPC)}.

  \item \textbf{Есть ли оптимизация скорости БПЛА.}
  В статье фокус на \emph{planning+tracking} и адаптацию MPC через RL, а не на явную постановку \emph{time-optimal} (минимизация времени) или отдельную оптимизацию профиля скорости «как можно быстрее до финиша».
  Практически скорость/динамика учитываются внутри ANMPC через ограничения и веса, но «минимизировать время прохождения траектории» как явная цель (в формате $T\to\min$) в первой постановке статьи не является основной заявленной задачей.

  \item \textbf{Есть ли компенсация аэродинамических помех.}
  В явном виде «аэродинамическая компенсация через модель аэродинамики/нейросеть» не выделена как отдельный модуль.
  Зато заявляется работа в сложной среде и робастность/устойчивость за счёт \emph{адаптивной} настройки MPC посредством RL.

  \item \textbf{Есть ли открытый датасет / как собирались данные.}
  Прямых указаний на открытый датасет/репозозиторий кода в тексте не обнаружено.
  Типичный способ воспроизведения — \textbf{генерировать данные в симуляторе}:
  состояния $x$, команды управления $u$, эталонные траектории, возмущения (ветер), ошибки слежения, затраты MPC; затем использовать их для обучения RL-агента, который подбирает веса MPC (или для imitation learning).
\end{itemize}

\section*{2) \textit{Aerodynamic Effects Compensation on Multi-Rotor UAVs Based on a Neural Network Control Allocation Approach} \href{https://www.ieee-jas.net/en/article/doi/10.1109/JAS.2021.1004266}{ref}}
\begin{itemize}
  \item \textbf{Какая задача решалась.}
  Компенсировать \textbf{аэродинамические эффекты} (в частности связанные с inflow / flapping и т.п.), которые обычно игнорируются в классической \emph{матрице микширования} (control allocation), из-за чего растёт ошибка аллокации и падает качество управления.

  \item \textbf{Использовалось ли ИИ.}
  Да: \textbf{нейросеть} используется как \emph{function fitting} для замены классической allocation/mixing matrix, чтобы по виртуальным управляющим воздействиям (тяга/моменты) выдавать корректные команды роторам с учётом аэродинамики. Обучение выполняется \textbf{оффлайн}.

  \item \textbf{Есть ли оптимизация скорости БПЛА.}
  Основной акцент не на time-optimal/скоростной оптимизации прохождения траектории, а на \textbf{улучшении качества управления/слежения} через уменьшение ошибки control allocation при аэродинамических эффектах.

  \item \textbf{Есть ли компенсация аэродинамических помех.}
  Да, это ядро работы: аэродинамика компенсируется \textbf{на уровне control allocation} — вместо упрощённой матрицы микширования используется нейросетевая аппроксимация инверсии «виртуальные усилия $\to$ скорости роторов» с учётом аэродинамических составляющих тяги/моментов.

  \item \textbf{Есть ли открытый датасет / как собирались данные.}
  В тексте явно говорится про \textbf{offline training}, но признаков публичного датасета/репозитория не видно.
  Чтобы повторить самостоятельно, тебе нужно собрать пары вида:
  \[
    (\text{состояния/скорости/ориентация},\ \text{вирт. тяга и моменты}) \mapsto (\omega_1,\dots,\omega_4),
  \]
  где целевые $\omega_i$ получаются либо из более полной аэродинамической модели (inflow+flapping), либо из стендовых/полетных логов с оценкой действительных сил/моментов.
\end{itemize}

\section*{3) \textit{Neural Network-PSO-based Velocity Control Algorithm for Landing UAVs on a Boat} \href{https://arxiv.org/pdf/2309.13679}{ref}}
\begin{itemize}
  \item \textbf{Какая задача решалась.}
  Обеспечить \textbf{быстрое и точное} управление \textbf{скоростью} для посадки БПЛА на \emph{движущуюся платформу} (лодку) в GPS-denied условиях, минимизируя ручную настройку регуляторов и опираясь на дешёвые датчики (камера + датчик высоты).

  \item \textbf{Использовалось ли ИИ.}
  Да: связка \textbf{PSO + Neural Network}.
  PSO оптимизирует параметры velocity-controller (переменные PID + параметры профиля снижения), а нейросеть учится \emph{интерполировать/предсказывать} параметры регулятора для непрерывного диапазона высот/скоростей.

  \item \textbf{Есть ли оптимизация скорости БПЛА.}
  Да, прямо: оптимизируется \textbf{velocity control} (скорости по осям) для \emph{быстрой} стыковки/посадки; приводится достижимая высокая скорость слежения/сближения (например, целевые скорости порядка нескольких м/с, вплоть до $\sim 8.9$ м/с в экспериментах).

  \item \textbf{Есть ли компенсация аэродинамических помех.}
  Специального аэродинамического компенсатора (типа drag/flapping model) нет, но учитываются реальные возмущения среды: упоминаются испытания и устойчивость даже при \textbf{ветровых условиях} (в контексте задач посадки на лодку).

  \item \textbf{Есть ли открытый датасет / как собирались данные.}
  Публичный датасет/код по явным признакам не приложен.
  Данные собирались через \textbf{sim-to-real}: оптимизация и обучение делаются в симуляции, затем валидируются полевыми тестами; тебе для воспроизведения нужно логировать:
  \[
    (\text{ошибка до цели по } x,y,z,\ \hat{v}_\text{цели},\ h,\ \text{ограничения}) \to (K_p,K_i,K_d,\alpha,\beta),
  \]
  а также фактические траектории/скорости/ошибки посадки.
\end{itemize}

\section*{4) \textit{A hierarchical ORCA framework for Multi-UAV navigation in unstructured environments with velocity optimization and local minima avoidance} \href{https://www.sciencedirect.com/science/article/pii/S0957417425028210}{ref}}
\begin{itemize}
  \item \textbf{Какая задача решалась.}
  Координированная навигация группы БПЛА без столкновений в \textbf{неструктурированных} средах с учётом \textbf{динамических ограничений}, избегание локальных минимумов, улучшение плавности движения.

  \item \textbf{Использовалось ли ИИ.}
  Нет, это алгоритмический/оптимизационный подход на базе ORCA (rule/optimization-based), без обучения модели на данных.

  \item \textbf{Есть ли оптимизация скорости БПЛА.}
  Да: вводится \textbf{VS-ORCA} (velocity-smoothed ORCA), где специально борются с «скачками скорости» (velocity mutation) и оптимизируют изменение скорости, чтобы получить более гладкие velocity inputs и улучшить показатели времени/длины/ошибок.

  \item \textbf{Есть ли компенсация аэродинамических помех.}
  Прямого аэродинамического компенсатора нет, но рассматриваются внешние возмущения среды: анализируются траектории/ускорения при разных \textbf{ветровых} условиях; показано, что алгоритм сохраняет корректность движения/избегание столкновений при ветре.

  \item \textbf{Есть ли открытый датасет / как собирались данные.}
  Это симуляционный алгоритмический paper; датасет в привычном ML-смысле обычно не требуется.
  Для воспроизведения тебе нужны \textbf{сцены/карты}, параметры препятствий/рельефа, динамические ограничения БПЛА, сценарии ветра, и логи: позиции, скорости, ускорения, коллизии/фейлы, время/длина пути.
\end{itemize}

\section*{5) \textit{Enhancing UAV Path Planning Efficiency through Adam-Optimized Deep Neural Networks for Area Coverage Missions} \href{https://www.sciencedirect.com/science/article/pii/S1877050924006707}{ref}}
\begin{itemize}
  \item \textbf{Какая задача решалась.}
  Генерация траекторий для миссий \textbf{покрытия области} (area coverage): сделать траектории более гладкими/точными, уменьшить резкие изменения направления и улучшить метрики ошибки координат (MSE/MAE/RMSE/$R^2$).

  \item \textbf{Использовалось ли ИИ.}
  Да: \textbf{DNN} для генерации траектории; обучение оптимизатором \textbf{Adam}; сравнение активаций (tanh/sigmoid/ReLU).

  \item \textbf{Есть ли оптимизация скорости БПЛА.}
  В явном виде оптимизации профиля скорости (как «минимум времени/максимум скорости при ограничении на отклонение») нет; оптимизируется точность/плавность траектории в координатах, а не time-optimal прохождение.

  \item \textbf{Есть ли компенсация аэродинамических помех.}
  Нет отдельного модуля компенсации аэродинамики; обсуждаются условия среды/препятствия на уровне постановки, но не как аэродинамический компенсатор.

  \item \textbf{Есть ли открытый датасет / как собирались данные.}
  Упоминается датасет сенсорных данных/траекторий как основа обучения, но без явной ссылки на публичную выгрузку.
  По тексту видно, что можно собрать самому: сгенерировать/записать пары «входные сенсорные/пространственные признаки $\to$ координаты траектории», разделить на train/test и обучать DNN.
\end{itemize}

\section*{6) \textit{Applying neural networks as direct controllers in position and trajectory tracking algorithms for holonomic UAVs} \href{https://pmc.ncbi.nlm.nih.gov/articles/PMC11993649/#Abs1}{ref}}
\begin{itemize}
  \item \textbf{Какая задача решалась.}
  Сравнить разные нейросети как \textbf{прямые контроллеры} (standalone control algorithms) для \textbf{позиционного} и \textbf{траекторного слежения} квадрокоптера (holonomic UAV), оценить точность и вычислительную эффективность.

  \item \textbf{Использовалось ли ИИ.}
  Да: тестируются различные архитектуры нейросетей как \textbf{контроллеры}. В тексте явно фигурирует MLP, обсуждаются эффекты размера скрытого слоя и проблемы обобщения.

  \item \textbf{Есть ли оптимизация скорости БПЛА.}
  Скорость учитывается через сценарии слежения за «tracked point» на разных скоростях (например, 0.5--3.5 м/с), и оценивается, где БПЛА не может догнать по скорости/радиусу. Но это скорее \textbf{оценка работоспособности контроллера на разных скоростях}, а не постановка time-optimal задачи «прилететь быстрее».

  \item \textbf{Есть ли компенсация аэродинамических помех.}
  Явного аэродинамического компенсатора (drag/flapping NN и т.п.) нет; речь про неопределённости/ошибки слежения и устойчивость нейроконтроллера.

  \item \textbf{Есть ли открытый датасет / как собирались данные.}
  Датасет формировался \textbf{симуляцией} на основе алгоритма позиционного слежения (artificial potential field) и разнообразия траекторий/скоростей; приведены детали: использовано 52 комбинации последовательностей, суммарно 116\,000 сэмплов. Публичную ссылку на датасет/код в явном виде по тексту не видно.
  Для ручного воспроизведения: генерируешь в симуляторе множество «tracked point» траекторий (прямоуг., треуг., круг, восьмёрка), скорости, относительные положения/ориентации; логируешь пары
  $(\Delta p, \Delta \psi, v_\text{ref}, \ldots)\to u$
  (управления) и обучаешь нейросеть как прямой контроллер.
\end{itemize}

%========================================================
% Важно (честно про файлы)
%========================================================
%\section*{Примечание по файлам}
%В папке с загруженными материалами доступны 6 PDF-файлов (перечислены выше).
%Файлов arXiv \texttt{2003.13801v1.pdf}, \texttt{2112.13396v1.pdf}, \texttt{2401.02899v1.pdf}, \texttt{2510.05698v1.pdf} среди реально загруженных в окружение сейчас нет, поэтому я не могу сделать по ним корректные памятки без повторной загрузки.
